\section{Derive addresses from the stored shape}
For zero-based indices, the byte addresses are
\begin{align*}
 \operatorname{addr}(A_{i,k})&=a+4(iK+k),\\
 \operatorname{addr}(B_{k,j})&=b+4(kN+j),\\
 \operatorname{addr}(C_{i,j})&=c+8(iN+j).
\end{align*}
Here \(a,b,c\) are full addresses supplied by the caller. The coefficient outside
each parenthesis is a storage width in bytes. It is unrelated to the tile edge or
the numerical value stored at that address.

For the worked example, suppose \(a=\mathtt{0x10010000}\),
\(b=\mathtt{0x10010480}\), and \(c=\mathtt{0x10010900}\).
The three addresses used for \(C_{1,0}\) in \(A\) are \(a+12,a+16,a+20\).
The matching addresses in \(B\) are \(b,b+8,b+16\). Its output address is
\(c+16\). The right-input column advances by \(4N=8\) bytes while the
left-input row advances by four bytes.

At function entry, load the descriptor fields once into working registers.
Compute a row offset with ordinary bounded integer arithmetic, then use a
64-bit address addition to add that offset to the supplied base. The dimensions
are small, but the base address may be above four gibibytes. A thirty-two-bit
addition can discard the high address bits even when the offset is only four.

\begin{table}[htbp]
\centering
\caption{Pointer movement while a baseline dot product keeps \(i,j\) fixed.}
\tagpdfsetup{table/header-rows={1}}
\begin{tabular}{lll}
\toprule
Pointer & Initial element & Change after one product\\
\midrule
Left input & \(A_{i,0}\) & Add 4 bytes\\
Right input & \(B_{0,j}\) & Add \(4N\) bytes\\
Output & \(C_{i,j}\) & No change until the sum is stored\\
\bottomrule
\end{tabular}
\end{table}

The debugger can display the same eight output bytes as two words or one
doubleword. For positive result 58, the low word is 58 and the high word is zero
in this little-endian environment. For result \(-1\), both words are
\texttt{ffffffff}. Inspect both halves when checking a result larger than a word.

\begin{MARSExercise}{Move the selected windows}
Move each supplied base past one unused element of its allocation. Recompute the
three starting addresses for \(C_{1,0}\). Which dimension determines the right
column stride? Which part of the calculation changes when only the output base
moves? Keep the inactive prefix unchanged.
\end{MARSExercise}
